% Page 052

\seite{52}

gelungen, den Frieden zu erhalten, da plötzlich hieß es: 'König Mobil!'\ob
der Kaiser rief, allzeit eilten zu den Waffen. Deutschland fühlte sich eins\ob
mit seinem Kaiser. Der Krieg an Frankreich war erklärt, der alte Erbfeind\ob
Deutschlands. Deutschlands Heer stand in wenigen Tagen an Rußlands Grenze.\ob
Der Durchbruch der Franzosen durch Belgien, wurde von den\ob
deutschen ferngehalten. Deutschland war gezwungen durch Belgien zu\ob
ziehen, da es mit Frankreich und England verbündet war. In der zweiten\ob
Woche des Monats August bekam unser Dorf viel Einquartierung, es\ob
waren sächsische Truppen, die nach Belgien zogen. Mit Pferden sollen\ob
11--1200 Mann im Dorfe. Hier aus dem Dorfe rückten 12--16 ein Reservisten.\ob
Während der Einquartierung taten die Leute alles auf, die Soldaten gut\ob
zu bewirten. Weil die Ernte noch nicht eingebracht war, um manche Familien\ob
durch die Einberufung die besten Arbeitskräfte verloren, halfen sich die Leute\ob
gegenseitig aus. 2--3 Familien arbeiteten zusammen mit Mähmaschinen und\ob
fuhren auch das Getreide zusammen ein, da die besten Arbeitspferde vom Heere\ob
angekauft waren. Manchmal halfen sogar während der Einquartierung die\ob
Soldaten selbst. Die Mägde und Knechte schon auch helfen. darum fiel der\ob
Schulunterricht 3 Wochen ganz aus, dennoch hatte Ostern frei, bis zum Beginn\ob
des Herbstferien. Am Samstag durfte Futter gearbeitet werden. Als so viele\ob
Züge mit Soldaten nach Osten und Westen nach dem Kriegsschauplatz fuhren,\ob
wurde im Dorf\unsicher{} gesammelt, man sammelte man im Dorfe für Fleisch, Speck,\ob
Brot, Leinen, Strümpfe und brachte die Sachen nach der Kirche zur Verteilung.\ob
Die Mädchen machten leinene Binden fürs rote Kreuz. Als zu Beginn des Win\ohb
ters für die Soldaten warme Kleidung beschafft werden mußte, regten sich\ob
viele fleißige Hände. Es wurde Geld gesammelt und gestrickt und genäht,\ob
Pulswärmer, Handschuhe, Leibbinden, Lungenwärmer, wollene Unterkleider, Strümpfe,\ob
Kniewärmer, Ohrenschützer. Jedes Kind wollte Patin sein des\unsicher\ob
ein reiches, schönes Weihnachtspaket von der Jugend, wobei die Jungfrau süße\ob
Zigaretten, Tabak etc. sorgten. Die Schulkinder sammelten für das Päckchen eifrig\ob
für Krieger. Zweimal am Tage rief die Glocke die Kinder zur Kapelle zum Beten\ob
und betete mit, selbst in der größten Arbeit und Hitze war die Kapelle gefüllt.\ob
In der Pfarrkirche wird fortwährend gesammelt und von Zeit zu Zeit erhalten\ob
die im Krieg der Pfarrei beteten. Während allgemein in der Zeit viel für das\ob
ganze Heer gebetet wird, ist die Teilnahme besonders groß, wenn das Seelenamt\ob
für einen Krieger gehalten wird. Zuerst aus der Pfarrei Gefallene: P{[eter]}
Gelesen      N.\unsicher, geschmückt mit dem eisernen Kreuz, gestorben, während bis jetzt im ganzen
19.2.15      Dekanat 10 aus der Pfarrei gefallen sind, 5 werden längere Zeit vermißt.\ob
LMG
\pend
